\Needspace{6\baselineskip}
\begin{center}\SourceLeaf{248}\pdfbookmark[2]{176.}{article-176}176.\end{center}
Iam perinde, ut positis valoribus incognitarum determinatis quibuscunque,
cuivis systemati valorum functionum $V$, $V'$, $V''$ etc. ante observationem
factam probabilitas determinata competit, ita vice versa, postquam ex
observationibus valores determinati functionum prodierunt, ad singula
systemata valorum incognitarum, e quibus illi demanare potuerunt, probabilitas
determinata redundabit: manifesto enim systemata ea pro magis probabilibus
habenda erunt, in quibus eventus eius qui prodiit exspectatio maior affuerat.
Huiusce probabilitatis aestimatio sequenti theoremati innititur:

\emph{Si posita hypothesi aliqua $H$ probabilitas alicuius eventus determinati
$E$ est $=h$, posita autem hypothesi alia $H'$ illam excludente et per se aeque
probabili eiusdem eventus probabilitas est $=h'$: tum dico, quando eventus $E$
revera apparuerit, probabilitatem, quod $H$ fuerit vera hypothesis, fore ad
probabilitatem, quod $H'$ fuerit hypothesis vera, ut $h$ ad $h'$}\HistoricalNoteEight{Handschriftliche Bemerkung:] Hätten die Hypothesen $H$, $H'$ an sich (d. i. vor dem Eintreten von $E$ oder vor erlangter Kenntniss von diesem Eintreten), ungleiche Wahrscheinlichkeiten $\mu$, $\mu'$ gehabt, so wird man ihnen, nach der Erscheinung von $E$, Wahrscheinlichkeiten beilegen müssen, die den Producten $\mu h$, $\mu' h'$ proportional sind.}.

Ad quod demonstrandum supponamus, per distinctionem omnium circumstantiarum,
a quibus pendet, num $H$ aut $H'$ aut alia hypothesis locum habeat, utrum eventus
$E$ an alius emergere debeat, formari systema quoddam casuum diversorum, qui
singuli per se (i. e. quamdiu incertum est, utrum eventus $E$ an alius
proditurus sit) tamquam aeque probabiles considerandi sint, hosque casus ita
distribui,
\begin{center}
\small
\begin{tabular}{c|c|c}
\shortstack{ut inter ipsos\\reperiantur} & \shortstack{ubi locum habere\\debet hypothesis} & \shortstack{cum modificationibus talibus,\\ut prodire debeat eventus} \\
\hline
$m$ & $H$ & $E$ \\
$n$ & $H$ & ab $E$ diversus \\
$m'$ & $H'$ & $E$ \\
$n'$ & $H'$ & ab $E$ diversus \\
$m''$ & ab $H$ et $H'$ diversa & $E$ \\
$n''$ & ab $H$ et $H'$ diversa & ab $E$ diversus
\end{tabular}
\end{center}
Tunc erit $h=\dfrac{m}{m+n}\mathpunct{\raisebox{.45ex}{,}}$
$h'=\dfrac{m'}{m'+n'}\mathpunct{\raisebox{.45ex}{;}}$ porro ante eventum
cognitum probabilitas \SourceLeaf{249}hypothesis $H$ erat
$=\dfrac{m+n}{m+n+m'+n'+m''+n''}\mathpunct{\raisebox{.45ex}{,}}$ post eventum
cognitum autem, ubi casus $n$, $n'$, $n''$ e possibilium numero abeunt, eiusdem
hypothesis probabilitas erit $=\dfrac{m}{m+m'+m''}\mathpunct{\raisebox{.45ex}{;}}$
perinde hypothesis $H'$ probabilitas ante et post eventum resp. exprimetur per
$\dfrac{m'+n'}{m+n+m'+n'+m''+n''}$ et $\dfrac{m'}{m+m'+m''}\mathpunct{:}$
quoniam itaque hypothesibus $H$ et $H'$ ante eventum cognitum eadem probabilitas
supponitur, erit $m+n=m'+n'$, unde theorematis veritas sponte colligitur.

Iam quatenus supponimus, praeter observationes $V=M$, $V'=M'$, $V''=M''$ etc.
nulla alia data ad incognitarum determinationem adesse, adeoque omnia systemata
valorum harum incognitarum ante illas observationes aeque probabilia fuisse,
manifesto probabilitas cuiusvis systematis determinati post illas observationes
ipsi $\Omega$ proportionalis erit. Hoc ita intelligendum est, probabilitatem,
quod valores incognitarum resp. iaceant inter limites infinite vicinos $p$ et
$p+\mathrm{d}p$, $q$ et $q+\mathrm{d}q$, $r$ et $r+\mathrm{d}r$, $s$ et
$s+\mathrm{d}s$ etc., exprimi per $\lambda\Omega\,\mathrm{d}p\,\mathrm{d}q\,\mathrm{d}r\,\mathrm{d}s$
etc., ubi $\lambda$ erit quantitas constans a $p$, $q$, $r$, $s$ etc. independens.
Et quidem manifesto erit $\dfrac{1}{\lambda}$ valor integralis ordinis
$\nu^{ti}$: $\int^{\nu}\Omega\,\mathrm{d}p\,\mathrm{d}q\,\mathrm{d}r\,\mathrm{d}s\, .\, .\, .\, .$,
singulis variabilibus $p$, $q$, $r$, $s$ etc. a valore $-\infty$ usque ad valorem
$+\infty$ extensis.
